%% explanation_v4_trial.tex — LiDAR-observation single-ownship environment.
% !TEX program = pdflatex

\documentclass[11pt, a4paper]{article}
\usepackage{amsmath, amssymb}
\usepackage{booktabs}
\usepackage{geometry}
\usepackage{microtype}
\usepackage{parskip}
\usepackage{graphicx}
\usepackage[hidelinks]{hyperref}
\usepackage{algorithm}
\usepackage{algpseudocode}
\geometry{margin=2.5cm}

\title{\textbf{v4\_trial Environment: Technical Specification}}
\author{}
\date{}

\begin{document}
\maketitle

% ─────────────────────────────────────────────────────────────────────────────
\section{Overview}
% ─────────────────────────────────────────────────────────────────────────────

The v4\_trial environment shares the same scenario setup as v3\_non\_changing:
a single ownship (slot~0) is controlled by the agent, all surrounding traffic
is uncontrolled, and the sector polygon, aircraft density, and spawning logic
are identical.
Two deliberate changes are made to the interface:

\begin{enumerate}
    \item \textbf{Observation.} The four explicit per-neighbour CPA features of
          v3\_non\_changing are replaced by a dual-channel LiDAR sensor:
          36 rays detecting aircraft proximity and 36 rays detecting the sector
          boundary, giving 77 features per frame.
          Temporal context is provided by stacking 4 consecutive frames,
          yielding a 308-dimensional input to the policy network.
    \item \textbf{Reward.} The safety penalty is now driven by the minimum
          LiDAR aircraft reading rather than a precomputed CPA distance.
          A navigation-progress term and a constant alive penalty replace the
          heading-drift penalty of v3\_non\_changing, following the reward
          structure of the \textsc{navigation-mappo-rl} reference implementation.
\end{enumerate}

The action space (Discrete(9) ATCO-style clearances) and scenario parameters
are unchanged.
The central research question is whether a policy can learn effective conflict
avoidance from a raw geometric proximity sensor alone, without access to
precomputed CPA distances or intruder velocity information in the observation.

% ─────────────────────────────────────────────────────────────────────────────
\section{Environment Loop}
% ─────────────────────────────────────────────────────────────────────────────

\begin{algorithm}[h]
\caption{\textsc{Reset} — initialise a new episode}
\begin{algorithmic}[1]
\State Generate random convex polygon (5--7 vertices)
\State area $\sim\mathcal{U}(10\,000,\,80\,000)$\,km$^2$,\quad
       density $\sim\mathcal{U}(5\,000,\,15\,000)$\,km$^2$/aircraft
\State $N \leftarrow \operatorname{clip}\!\bigl(\lfloor\text{area}/\text{density}\rceil,\;2,\;15\bigr)$
\State Compute $T_\text{max}$ from sector diameter and cruise speed
\For{each slot $s = 1,\ldots,N$}
    \State Spawn aircraft on perimeter arc $s$ with uniform jitter
\EndFor
\State Ownship $\leftarrow$ aircraft in slot~0 \quad\textit{(fixed for the episode)}
\State Compute initial LiDAR readings $(\boldsymbol{\ell}^\text{ac}_0,\,\boldsymbol{\ell}^\text{bnd}_0)$
\State \Return observation $\mathbf{o}(\text{ownship})$
\end{algorithmic}
\end{algorithm}

\begin{algorithm}[h]
\caption{\textsc{Step}(action) — advance one RL step (30\,s simulated)}
\begin{algorithmic}[1]
\State Spawn any pending replacement aircraft
\If{action is heading offset $\delta$}
    \State $\psi_\text{cmd} \leftarrow \psi_\text{dest}(\text{ownship}) + \delta$
\ElsIf{action is Mach adjustment $\Delta M$}
    \State $M_\text{cmd} \leftarrow \operatorname{clip}(M_\text{cmd} + \Delta M,\;0.70,\;0.82)$
\EndIf
\State \textit{(Traffic aircraft receive no commands)}
\State Advance BlueSky simulation $60\times0.5$\,s $= 30$\,s
\For{each aircraft that exited the sector}
    \If{exiting aircraft is ownship}
        \State Apply exit penalty if $|\Delta\psi_\text{dest}| > 90^\circ$
    \EndIf
    \State Schedule replacement
\EndFor
\State Update ownship callsign if slot~0 was replaced
\State Compute LiDAR readings $(\boldsymbol{\ell}^\text{ac},\,\boldsymbol{\ell}^\text{bnd})$ around updated ownship
\State $r \leftarrow r_\text{nav} + r_\text{safe} + r_\text{alive} + r_\text{work} + r_\text{exit}$
\State \Return $\mathbf{o}(\text{ownship}),\; r,\;
               \text{truncated}=(\text{step\_count} \geq T_\text{max})$
\end{algorithmic}
\end{algorithm}

\paragraph{Step duration.}
Each RL step covers $60 \times 0.5\,\text{s} = 30$\,s of simulated time.
This ensures that a commanded 30° heading change is physically complete before
the reward is evaluated, so the policy observes the full geometric consequence
of each instruction.

% ─────────────────────────────────────────────────────────────────────────────
\section{Action Space}
% ─────────────────────────────────────────────────────────────────────────────

Identical to v3\_non\_changing.
Nine discrete ATCO-style clearances applied to the ownship only:
\begin{equation}
    \delta \in \{-30^\circ,\,-15^\circ,\,0^\circ,\,+15^\circ,\,+30^\circ,\,\varnothing\},
    \qquad
    \Delta M \in \{-0.04,\,-0.02,\,+0.02\}.
\end{equation}
Workload costs: 1.0 heading turns, 0.5 speed adjustments, 0 direct and hold.

% ─────────────────────────────────────────────────────────────────────────────
\section{LiDAR Sensor}
% ─────────────────────────────────────────────────────────────────────────────

The LiDAR casts $N_\text{rays} = 36$ rays uniformly around the ownship in its
ego-centric frame (ray~0 = dead ahead; rays proceed clockwise at $10^\circ$
intervals).
Two channels are computed per ray, mirroring the agent and boundary channels
of \textsc{navigation-mappo-rl}.

\paragraph{Aircraft channel $\boldsymbol{\ell}^\text{ac}$.}
Each traffic aircraft is modelled as a circle of radius $r_\text{det} = 5$\,NM.
For a ray in direction $(\sin\theta,\cos\theta)$ in the ego frame and an
intruder at ego-frame position $(p_x, p_y)$, the hit distance is
\begin{equation}
    t_\text{hit} = \max\!\bigl(0,\; t - \sqrt{r_\text{det}^2 - p_\perp^2}\bigr),
    \quad
    t = p_x\sin\theta + p_y\cos\theta,
    \quad
    p_\perp^2 = p_x^2 + p_y^2 - t^2,
\end{equation}
whenever $t > 0$ and $p_\perp^2 < r_\text{det}^2$.
The reading for ray $r$ is $\ell^\text{ac}_r = \min_k(t_{\text{hit},k})/d_\text{max}$,
clamped to $[0,1]$ and taking the minimum over all traffic aircraft~$k$.
A value of 1 indicates no aircraft within $d_\text{max} = 30$\,NM; 0 indicates
contact.

\paragraph{Boundary channel $\boldsymbol{\ell}^\text{bnd}$.}
The same 36 rays are cast against the sector polygon boundary using ray–polygon
intersection.
Each reading is $\ell^\text{bnd}_r = d_\text{boundary}/d_\text{max}$, clamped
to $[0,1]$, where $d_\text{boundary}$ is the distance to the nearest polygon
edge along ray $r$.
A value of 0 indicates the ownship is at the boundary; 1 indicates the
boundary is beyond $d_\text{max}$ or no intersection exists within range.

% ─────────────────────────────────────────────────────────────────────────────
\section{Observation Space}
% ─────────────────────────────────────────────────────────────────────────────

\paragraph{Single frame (77 floats).}

\begin{table}[h]
\centering
\caption{Per-frame observation vector (77 features).}
\label{tab:obs}
\smallskip
\begin{tabular}{clll}
\toprule
\textbf{Index} & \textbf{Feature} & \textbf{Definition} & \textbf{Range} \\
\midrule
\multicolumn{4}{l}{\textit{Ownship state (5)}} \\[2pt]
0 & $\sin\Delta\psi$ & $\sin(\psi_\text{dest}-\psi_\text{cmd})$ & $[-1,1]$ \\[2pt]
1 & $\cos\Delta\psi$ & $\cos(\psi_\text{dest}-\psi_\text{cmd})$ & $[-1,1]$ \\[2pt]
2 & Speed ratio      & $M_\text{cmd}/M_\text{nominal}$           & $[0,\sim\!1]$ \\[2pt]
3 & Turn progress    & $\operatorname{clip}((\psi_\text{cmd}-\psi_\text{actual})/30^\circ,\,-1,1)$ & $[-1,1]$ \\[2pt]
4 & Time to exit     & $\min(d_\text{boundary}/(v\cdot T_\text{look}),\,1)$ & $[0,1]$ \\[4pt]
\multicolumn{4}{l}{\textit{Aircraft LiDAR — 36 ego-centric rays (ray 0 = ahead, clockwise):}} \\[2pt]
$5\ldots40$   & $\ell^\text{ac}_0,\ldots,\ell^\text{ac}_{35}$   & Distance to nearest aircraft & $[0,1]$ \\[4pt]
\multicolumn{4}{l}{\textit{Boundary LiDAR — same 36 rays against sector polygon:}} \\[2pt]
$41\ldots76$  & $\ell^\text{bnd}_0,\ldots,\ell^\text{bnd}_{35}$ & Distance to sector boundary  & $[0,1]$ \\
\bottomrule
\end{tabular}
\end{table}

\paragraph{Temporal context — frame stacking.}
Following the \textsc{navigation-mappo-rl} approach (\texttt{history\_length}~$= 4$),
four consecutive frames are stacked by \texttt{VecFrameStack(4)} in the
training script, giving a 308-dimensional observation fed to the policy:
\begin{equation}
    \mathbf{o}_t^\text{stacked}
    = \bigl[\mathbf{o}_{t},\;\mathbf{o}_{t-1},\;\mathbf{o}_{t-2},\;\mathbf{o}_{t-3}\bigr]
    \in \mathbb{R}^{308}.
\end{equation}
This allows the policy to infer intruder velocities from consecutive LiDAR
snapshots without explicit velocity features in the observation.

% ─────────────────────────────────────────────────────────────────────────────
\section{Reward Function}
% ─────────────────────────────────────────────────────────────────────────────

\begin{equation}
    r_t = r_\text{nav} + r_\text{safe} + r_\text{alive} + r_\text{work} + r_\text{exit}.
\end{equation}

\paragraph{Navigation term.}
Adapted from the progress reward of \textsc{navigation-mappo-rl}:
\begin{equation}
    r_\text{nav}
    = w_\text{nav}\cdot\cos\Delta\psi_\text{actual}\cdot\frac{M_\text{cmd}}{M_\text{nominal}},
\end{equation}
where $\Delta\psi_\text{actual} = \psi_\text{dest} - \psi_\text{actual}$.
The term is positive when the ownship flies toward its destination at nominal
speed and negative when flying away, providing a continuous navigation signal.
Scaling by the speed ratio encourages the agent to maintain cruise speed.

\paragraph{Safety term.}
Driven by the minimum aircraft LiDAR reading, converted back to nautical miles:
\begin{equation}
    d_\text{ac}^\text{min} = \min_r(\ell^\text{ac}_r)\cdot d_\text{max},
    \qquad
    r_\text{safe} = -w_\text{safe}\cdot
        \exp\!\left(-\frac{d_\text{ac}^\text{min}}{d_\text{sep}}\right).
\end{equation}
This uses the same exponential form as v3\_non\_changing but is evaluated on
the current geometric proximity rather than the predicted CPA distance.

\paragraph{Alive penalty.}
A small constant penalty per step, following \textsc{navigation-mappo-rl}:
\begin{equation}
    r_\text{alive} = -w_\text{alive}.
\end{equation}
This discourages the agent from remaining stationary and encourages efficient
sector transit.

\paragraph{Workload and exit terms.}
$r_\text{work} = -w_\text{work}\cdot c$, where $c$ is the discrete action
cost (identical to v3\_non\_changing).
$r_\text{exit} = -w_\text{exit}$ when the ownship exits the sector in the
wrong direction; zero otherwise.

\begin{table}[h]
\centering
\caption{Reward weights.}
\smallskip
\begin{tabular}{llr}
\toprule
$w_\text{safe}$  & Safety: LiDAR-proximity penalty           & $1.00$ \\
$w_\text{exit}$  & Wrong-direction exit (sparse)             & $0.50$ \\
$w_\text{nav}$   & Navigation: heading alignment $\times$ speed & $0.20$ \\
$w_\text{alive}$ & Alive cost per step                       & $0.05$ \\
$w_\text{work}$  & Workload: instruction cost                & $0.05$ \\
\bottomrule
\end{tabular}
\end{table}

% ─────────────────────────────────────────────────────────────────────────────
\section{Comparison with v3\_non\_changing}
% ─────────────────────────────────────────────────────────────────────────────

\begin{table}[h]
\centering
\caption{Key differences between v3\_non\_changing and v4\_trial.}
\smallskip
\begin{tabular}{lll}
\toprule
& \textbf{v3\_non\_changing} & \textbf{v4\_trial} \\
\midrule
Action space         & Discrete(9)                      & Discrete(9) \\
Step duration        & 30\,s                            & 30\,s \\
Obs per frame        & 21                               & 77 \\
Obs (stacked)        & 21                               & 308 (4 frames) \\
Obs [0:5]            & 5 ego features                   & 5 ego features \\
Obs [5:21]           & 4 neighbours $\times$ 4 CPA features & 36 aircraft LiDAR rays \\
Obs [21:77]          & —                                & 36 boundary LiDAR rays \\
CPA in obs           & Yes (explicit)                   & No \\
CPA in reward        & Yes                              & No \\
Safety reward        & $-w\cdot\exp(-d_\text{CPA}/d_\text{sep})$  & $-w\cdot\exp(-d_\text{lidar}/d_\text{sep})$ \\
Navigation reward    & $-w_\text{eff}\cdot(1-\cos\Delta\psi_\text{cmd})/2$ & $+w_\text{nav}\cdot\cos\Delta\psi_\text{actual}\cdot\text{speed\_ratio}$ \\
Alive penalty        & No                               & $-w_\text{alive}$ per step \\
Urgency matrix       & Yes (visualisation)              & Removed \\
\bottomrule
\end{tabular}
\end{table}

The central research question of v4\_trial is whether a policy can learn
effective conflict avoidance from a LiDAR proximity sensor alone.
Without explicit CPA features, the agent must infer the trajectory of
surrounding traffic from the spatial pattern of aircraft LiDAR readings across
consecutive stacked frames.
The boundary LiDAR channel gives the agent directional awareness of the sector
boundary in all directions, replacing the scalar time-to-exit computation used
in v3\_non\_changing with a richer spatial representation.
All other design decisions are held constant to isolate the effect of the
observation representation.

% ─────────────────────────────────────────────────────────────────────────────
\section{Training}
% ─────────────────────────────────────────────────────────────────────────────

Training uses PPO with four parallel environments (\texttt{SubprocVecEnv}),
network $[128,\,128]$, entropy coefficient $\beta_H = 0.02$, and all other
hyperparameters identical to v3\_non\_changing.
\texttt{VecFrameStack(4)} is applied after \texttt{VecNormalize}, stacking
four observation frames before feeding them to the policy.
Results are saved to \texttt{Runs\_saved/v4\_trial/}.

\end{document}
